\documentclass[11pt]{article}
\usepackage{amsmath,amssymb,amsthm,mathtools,hyperref}

\begin{document}
\section*{Three-dimensional tissue reconstruction}
\label{prob:077}
\label{rank:132}
\par\smallskip

\noindent\textbf{Domain.} Three-dimensional tissue inference \quad \textbf{Retained status/source context.} Open deformation-aware reconstruction and uncertainty problem; section alignment already has effective algorithms.
\paragraph{Definitions and assumptions.} A tissue occupies $\Omega\subset\mathbb{R}^3$ with molecular field $X:\Omega\to\mathbb{R}^p$. Section $i$ is sampled at measured depth $z_i$, deformed by an unknown in-plane map $\phi_i$, and observed with cell loss and assay noise. Maps belong to a bounded smooth diffeomorphism class with fixed orientation and landmark constraints. Tissue thickness, section spacing and missing sections are recorded. Adjacent sections contain distinct sampled cells, so cell identity is not a necessary one-to-one correspondence.
\paragraph{Problem statement.} Recover a field and deformations with controlled spatial and molecular errors from
\[Y_i(s)=X(\phi_i(s),z_i)+\eta_i(s),\qquad d_{\mathcal G}(\widehat X,X)\le\varepsilon.\]
Here $\mathcal G$ denotes the permitted unavoidable global coordinate transformations; the quotient identifies reconstructions related by those transformations. Specify how landmark density, section spacing and smoothness constrain uncertainty, and test reconstruction against independently imaged three-dimensional structures or withheld sections.
\paragraph{Additional binding assumptions and resolution criterion.} Define the quotient discrepancy by $d_{\mathcal G}(\widehat X,X):=\inf_{g\in\mathcal G}\|\widehat X-X\circ g\|_{L^2(\Omega)}$, with the permitted action and volume convention fixed. The action must make $X\circ g$ well-defined on the comparison domain; no symmetry or triangle-inequality property of this discrepancy is assumed without justification. Depth interpolation needs a quantitative smoothness bound; in-plane landmarks alone cannot identify tissue changes between absent sections. Distinguish uncertain deformations from genuine section-to-section biological variation.

\paragraph{Concrete-instance declarations.} In the displayed point-evaluation model, let $D_i$ denote the section coordinate domain and let $A_i\subseteq D_i$ denote the observed sampling locations after cell loss. The map $\phi_i$ takes section coordinates to in-plane tissue coordinates, with $(\phi_i(s),z_i)\in\Omega$ wherever the model is evaluated; observations are specified only for $s\in A_i$. This equation does not itself describe averaging over finite assay footprints or section thickness. Declare the sampling and loss mechanism, noise assumptions and treatment of depth uncertainty, without presuming independence or expression-independent loss. Declare quantitative field and deformation classes, including domains, regularity bounds, boundary conventions, landmark constraints and landmark uncertainty, sufficient to interpret the evaluations and discrepancy. The permitted coordinate transformations, their action on fields and deformations, their compatibility with depths and landmarks, and the comparison measure remain instance declarations. Before fitting or testing, specify the meaning and scope of the error guarantee and retain controlled spatial and molecular errors as joint requirements, with separate geometry, deformation, molecular and cell-correspondence scores under a common coordinate convention and predeclared scales and tolerances. Cell-correspondence scoring must not presume identical cells in adjacent sections. State which outcomes independent imaging or withheld sections measure, include reference uncertainty, and distinguish validation on observed planes from claims about unobserved volume. These declarations do not select a noise law, regularity class, symmetry action, guarantee interpretation or validation protocol that the supplied statement has left unspecified.

\paragraph{Alternative formulations.} (i) Use rigid or affine alignment for minimally distorted sections and reconstruct only anatomically defined regions. (ii) Use nonlinear registration with expression and histology, including explicit uncertainty for tears, folds and missing material. Interpolation across unobserved tissue needs a spatial regularity assumption; it cannot recover arbitrary hidden microstructure.
\paragraph{Obstructions and partial progress.} PASTE provides probabilistic alignment and integration, while the July 2026 review covers three-dimensional reconstruction as a continuing methodological area. Matching molecular patterns can confuse repeated anatomical structures or symmetries. Flexible deformation can achieve excellent apparent fit while distorting true geometry. A finite set of planes cannot determine arbitrary variation between planes. The narrowed target is resolution-limited reconstruction with landmarks and regularization, with cell correspondence, tissue geometry and molecular interpolation scored separately.
\paragraph{Sources (supplied provenance; not all rechecked).}
\begin{itemize}
\item Zhou et al., Computational analysis in spatial transcriptomics: methods and perspectives (July 2026). \url{https://pubmed.ncbi.nlm.nih.gov/42555503/} (Supplied recent review; explicitly covers these computational tasks and remaining challenges.)
\item Zeira et al., Alignment and integration of spatial transcriptomics data (2022). \url{https://www.nature.com/articles/s41592-022-01459-6} (PASTE establishes partial alignment progress.)
\end{itemize}


\section*{Binding catalogue conventions}
This edition revises both supplied catalogues and preserves original problem numbers 1--330. It adds one globally unique label \verb|\label{rank:x}| per problem, where smaller $x$ means greater provisional difficulty. A rank label is an anchor at that problem's section: \verb|\ref{rank:x}| returns the section number, not the literal difficulty number. The visible \texttt{rank:x} text and the index give the difficulty number itself.

The entries are \emph{parameterized mathematical research schemes}, including empirical prediction, causal identification, historical reconstruction, theory and engineering. They are not 330 certified unsolved mathematical conjectures. A rigorous instance fixes the objects and restrictions below; changing those objects changes the question. Established identities, finite calculations, non-identification constructions and solved special cases are not being posed anew as open. Some residual tasks concern informative bounds, biological realization or external validation rather than a new theorem.

Every main definition and problem statement received a mathematical editorial review. Each entry now has an additional binding-conditions paragraph; direct corrections are separately recorded in the accompanying audit. The alternative formulations inherit the applicable conventions, but are different targets unless equivalence is proved. Where an old broad phrase conflicts with an explicit restriction below it, the explicit restriction controls.

\textbf{Literature boundary.} Thirteen selective primary-literature checks cover eighteen entries in this revision; their access levels and scope are recorded explicitly. The supplied bibliography and its earlier source-access descriptions are retained as provenance, not represented as newly verified. Neither this review nor a source's description proves the absence of later solutions. The date of this edition is 13 October 2026; there is no claim of exhaustive current-status certification for all 330 entries.

\subsection*{Binding conventions for a concrete instance}
\paragraph{Objects, quantifiers and non-vacuity.}
An instance declares measurable state, input, output and parameter spaces; a fixed candidate family; admissible initial laws, interventions and environments; a finite horizon or specified asymptotic regime; observation and sampling laws; and performance metrics, positive normalization scales and positive tolerances. These are chosen before fitting or testing, not adjusted after seeing failures. Shared parameters must genuinely be shared across the declared validation domain. Finite dimension alone does not restrict arbitrary real-number encodings: specify a finite library, bounded description length, or regularity/complexity bounds. Exclude constant or no-signal solutions through a fixed informative-output or baseline-improvement requirement where appropriate. A universally quantified result cannot be established by choosing one trivial favorable instance.

\paragraph{Probability laws and data.}
Every expectation or probability refers to the specified stochastic law, including initial uncertainty, process noise, observation noise and sampling as appropriate. A deterministic equation is not a probability model without a declared random input or initial law. ``Observed'' targets in population-level formulas mean population observables $g_j^*$, not noisy estimates asserted to be exact. Let $\widetilde g_j$ be independent validation estimates and suppose
\[
\mathbb{P}\{\lvert\widetilde g_j-g_j^*\rvert\le r_j\text{ for every }j\}\ge1-\alpha.
\]
Then an uncertainty-aware sufficient certification rule for positive target tolerances $\varepsilon_j$ is
\[
\max_j\frac{|\widehat g_j-\widetilde g_j|+r_j}{\varepsilon_j}\le1.
\]
On that simultaneous confidence event it implies every population error is at most its tolerance, by the triangle inequality. Vector/norm versions use corresponding valid simultaneous bounds. Training uncertainty must also be included when the claim averages over repeated training data. Evaluation on finitely many conditions does not establish a continuous-domain supremum without covering and regularity bounds.

\paragraph{Dynamics and well-posedness.}
For ODEs and SDEs give coefficient domains, initial laws and existence/uniqueness conditions in the intended solution sense, or keep the entire admitted solution set. For stochastic jump models specify nonnegative off-diagonal rates, diagonal normalization, rate orientation and non-explosion. For PDEs specify spatial domains, boundary/initial conditions, conservation or source terms, and a weak or classical solution class. Moving-domain concentration equations include dilution/transport terms. A stationary law is used only for a time-homogeneous augmented model under existence assumptions; a stationary approximation to a driven system must be justified. Absorbing extinction states are retained. Survival-conditioned, quasi-stationary and unconditional laws are distinct and must not be substituted for one another.

\paragraph{Identification and causal meaning.}
Two parameter values are observationally equivalent when they induce the same full distribution under all allowed observation/intervention designs, not merely the same likelihood value for one sample. A target is identified only if constant on the equivalence class. Gauge symmetries, label switching and reference conventions are quotiented or fixed. Finite-data compatible sets use a declared confidence, likelihood or discrepancy rule; an empty set signals model rejection. Bayesian credible sets are not automatically confidence sets. A prior or regularizer does not create data identification. Bounds intended to resolve a question must be informative at a predeclared width or decision threshold.

A causal contrast uses well-defined intervention versions and a stated causal model. Random assignment, or explicit consistency, overlap and exchangeability assumptions, is required for causal interpretation from observations. Interference and time-varying confounding are represented when present. Fixing a mediator by intervention differs from conditioning on its observed value. Marginal potential-outcome laws do not determine individual cross-world effects or an arbitrary joint coupling. Ethically admissible designs, rather than unconstrained interventions, define the experimental menu.

\paragraph{Metrics, units, optimization and limits.}
Specify the sampling or reference measure for densities and integrals, the log base, and dimensionless arguments of logarithms. Total variation is $\sup_A|P(A)-Q(A)|$; Wasserstein distances require the declared ground metric and appropriate moment conditions; KL divergence uses its stated absolute-continuity convention. All normed readouts are scaled consistently, and error probabilities use $0<\alpha<1$. A ratio has a positive denominator or an explicit undefined/zero convention. Thermodynamic entropy production is nonnegative in expectation under the relevant hypotheses; individual stochastic realizations need not be. An infimum is not an attained minimum without compactness/coercivity and continuity or another existence argument. Fix the order of large-population, long-time and vanishing-mutation limits. Stability claims declare perturbation space and symmetries.

\subsection*{How the provisional difficulty order was assigned}
Four editorial components are scored from 1 (lower burden) to 5 (higher burden): (a) formal or inferential difficulty, (b) access to evidence and identification burden, (c) mechanistic/computational integration, and (d) independent validation and transfer. Their unweighted total orders the entries. For selected entries in the highest bands, an explicit editorial priority breaks equal totals; all other ties are broken by original problem number. Equal totals are \emph{not} evidence of reliably different true difficulties. Scores are judgments, not measurements, completion-time estimates or consensus rankings.

The object ranked is an informative, nondegenerate resolution of the \emph{scoped} scientific target, with the stated budget and validation domain. A sharp useful impossibility result can qualify; an already known identity, vacuous bound, arbitrary favorable model or lookup table cannot. Rank 330 is merely the least difficult relative assessment in this catalogue, not a claim of an easy biological solution. Broad historical titles do not automatically outrank precisely formulated mechanistic tasks.



\end{document}
